\section{Method}
\label{method}
This section describes the algorithmic details of our method.
To set up notation, consider a standard $K$-layer neural network with parameters~$\theta_k$ for layer~$k$. The stacked parameter vector $\vtheta=(\theta_1,\dots,\theta_K)$ specifies the entire model for a classification task with loss function~$l_m(x, y; \vtheta)$ on the test sample $(x,y)$. We call this the \emph{main task}, as indicated by the subscript of the loss function. 

We assume to have training data $(x_1,y_1),\dots,(x_n, y_n)$ drawn i.i.d.~from a distribution~$P$. Standard empirical risk minimization solves the optimization problem:
\begin{equation} \label{optimize}
\min_\vtheta \frac{1}{n}\sum_{i=1}^{n} l_m(x_i, y_i; \vtheta).
\end{equation}
Our method requires a \emph{self-supervised auxiliary task} with loss function~$l_s(x)$. 
In this paper, we choose the rotation prediction task \citep{gidaris2018unsupervised}, which has been demonstrated to be simple and effective at feature learning for convolutional neural networks.
The task simply rotates $x$ in the image plane by one of 0, 90, 180 and 270 degrees and have the model predict the angle of rotation as a four-way classification problem. 
Other self-supervised tasks in \Cref{related} might also be used for our method.

The auxiliary task shares some of the model parameters
$\vtheta_e=(\theta_1,\dots,\theta_\kappa)$ up to a
certain~$\kappa\in\{1,\dots,K\}.$
We designate those $\kappa$ layers as a
\emph{shared feature extractor}. 
The auxiliary task uses its own task-specific parameters 
$\vtheta_s = (\theta'_{\kappa+1},\dots, \theta'_{K})$. 
We call the unshared parameters $\vtheta_s$ the \emph{self-supervised task branch}, and $\vtheta_m=(\theta_{\kappa+1},\dots,\theta_K)$ the
\emph{main task branch}.
Pictorially, the joint architecture is a $Y$-structure with a shared bottom and two branches.
For our experiments, the self-supervised task branch has the same architecture as the main branch, except for the output dimensionality of the last layer due to the different number of classes in the two tasks.

Training is done in the fashion of multi-task learning \citep{caruana1997multitask}; the model is trained on both tasks on the same data drawn from~$P$. 
Losses for both tasks are added together, and gradients are taken for the collection of all parameters. The joint training problem is therefore
\begin{equation}
\label{optimize_train}
\min_{\vtheta_e,\vtheta_m,\vtheta_s}
\frac{1}{n}\sum_{i=1}^{n} 
l_m(x_i, y_i; \vtheta_m, \vtheta_e)
+ l_s(x_i; \vtheta_s, \vtheta_e).
\end{equation}

Now we describe the standard version of Test-Time Training on a single test sample $x$.
Simply put, Test-Time Training fine-tunes the shared feature extractor $\vtheta_e$ by minimizing the auxiliary task loss on $x$. 
This can be formulated as
\begin{equation}
    \label{optimize_test}
    \min_{\vtheta_e}l_s(x; \vtheta_s, \vtheta_e).
\end{equation}
Denote $\vtheta_e\opt$ the (approximate) minimizer of \autoref{optimize_test}.
The model then makes a prediction using the updated parameters 
$\vtheta(x) = (\vtheta_e\opt, \vtheta_m)$.
Empirically, the difference is negligible between minimizing \autoref{optimize_test} over $\vtheta_e$ versus over both $\vtheta_e$ and $\vtheta_s$.
Theoretically, the difference exists only when optimization is done with more than one gradient step.

Test-Time Training naturally benefits from standard data augmentation techniques. On each test sample $x$, we perform the exact same set of random transformations as for data augmentation during training, to form a batch only containing these augmented copies of $x$ for Test-Time Training.

\paragraph{Online Test-Time Training.}
In the standard version of our method, the optimization problem in
\autoref{optimize_test} is always initialized with parameters 
$\vtheta = (\vtheta_e,\vtheta_s)$ 
obtained by minimizing \autoref{optimize_train}.
After making a prediction on $x$, $\vtheta_e\opt$ is discarded.
Outside of the standard supervised learning setting, 
when the test samples arrive online sequentially, 
the online version solves the same optimization problem as in \autoref{optimize_test} to update the shared feature extractor $\vtheta_e$.
However, on test sample $x_t$, 
$\vtheta$ is instead initialized with $\vtheta(x_{t-1})$ updated on the previous sample $x_{t-1}$.
This allows $\vtheta(x_t)$ to take advantage of the distributional information available in $x_1, \dots,x_{t-1}$ as well as $x_t$.
